\chapter{1977年内蒙古卷（文）}

\begin{problem}
\begin{enumerate}
\item 计算 \(\sin\left(-990^\circ\right)\).
\item 化简 \(\sqrt{4-12a+9a^2}\).
\item 计算 \(\frac{(-27)^3\sqrt{\left(\frac{1}{16}\right)^{-\frac{1}{2}}}}{(-9)^0\left(\frac{3}{2}\right)^7\left(\frac{1}{2}\right)^{-5}}\).
\item 已知 \(\lg2=0.3010\)，\(f\left(x\right)=\log_{\sqrt{10}}x\)，求 \(f\left(5\right)\).
\item 求函数 \(y=\frac{\sqrt{1-\frac{2}{3}x}}{\sqrt[3]{3x-1}}\) 的定义域.
\item 已知 \(\sin\theta=0.6\)，求 \(\left(\sin\frac{\theta}{2}-\cos\frac{\theta}{2}\right)^2\) 的值.
\end{enumerate}
\end{problem}

\begin{solution}
\begin{enumerate}
\item 利用正弦函数的周期性，\(-990^\circ+3\times360^\circ=90^\circ\)，所以
\(\sin\left(-990^\circ\right)=\sin90^\circ=1\).

\item 先因式分解被开方数，得
\(4-12a+9a^2=\left(2-3a\right)^2\)，因此
\(\sqrt{4-12a+9a^2}=\left|2-3a\right|\).
再根据 \(2-3a\) 的符号分类讨论：当 \(a<\frac{2}{3}\) 时，\(2-3a>0\)；当 \(a=\frac{2}{3}\) 时，\(2-3a=0\)；当 \(a>\frac{2}{3}\) 时，\(2-3a<0\). 所以
\[
\sqrt{4-12a+9a^2}=
\begin{cases}
2-3a, & a<\frac{2}{3},\\
0, & a=\frac{2}{3},\\
3a-2, & a>\frac{2}{3}.
\end{cases}
\]

\item 因为 \((-9)^0=1\)，且
\(\left(\frac{1}{16}\right)^{-\frac{1}{2}}=\sqrt{16}=4\)，所以根式的值为 \(\sqrt{4}=2\). 另外，
\[
\begin{gathered}
\frac{(-27)^3\sqrt{\left(\frac{1}{16}\right)^{-\frac{1}{2}}}}{(-9)^0\left(\frac{3}{2}\right)^7\left(\frac{1}{2}\right)^{-5}}
=\frac{(-3)^9\cdot2}{3^7\left(\frac{1}{2}\right)^{7-5}}
=\frac{-3^9\cdot2}{3^7\left(\frac{1}{2}\right)^2}\\
=-3^2\cdot2\cdot4=-3^2\cdot2^3=-72.
\end{gathered}
\]

\item 由于 \(\sqrt{10}>0\) 且 \(\sqrt{10}\ne1\)，可用换底公式计算：
\[
\begin{gathered}
f\left(5\right)=\log_{\sqrt{10}}5
=\frac{\lg5}{\lg\sqrt{10}}
=\frac{\lg\frac{10}{2}}{\frac{1}{2}\lg10}\\
=2\left(\lg10-\lg2\right)
=2\left(1-0.3010\right)
=1.398
\end{gathered}
\]
其中 \(\lg10=1\)，故 \(f\left(5\right)=1.398\).

\item 分子中的平方根有意义，必须满足 \(1-\frac{2}{3}x\ge0\)，解得 \(x\le\frac{3}{2}\). 分母中的立方根对任意实数都有意义，但分母不能为零，因此还需满足 \(3x-1\ne0\)，即 \(x\ne\frac{1}{3}\). 两个条件取交集，函数的定义域为
\(\left(-\infty,\frac{1}{3}\right)\cup\left(\frac{1}{3},\frac{3}{2}\right]\).

\item 令 \(u=\frac{\theta}{2}\)，则
\[
\begin{gathered}
\left(\sin\frac{\theta}{2}-\cos\frac{\theta}{2}\right)^2
=\sin^2\frac{\theta}{2}+\cos^2\frac{\theta}{2}-2\sin\frac{\theta}{2}\cos\frac{\theta}{2}\\
=1-\sin\theta=1-0.6=0.4.
\end{gathered}
\]
\end{enumerate}
\end{solution}

\begin{problem}
\begin{enumerate}
\item 已知：梯形 \(ABCD\)，\(AD\parallel BC\)，对角线 \(AC\) 和 \(BD\) 相交于 \(E\)，过 \(E\) 作平行于底的直线，交 \(AB\) 于 \(M\)，交 \(CD\) 于 \(N\)，求证：\(ME=EN\).
\item 圆台上底直径为 \(12\mathrm{~cm}\)，下底直径为 \(24\mathrm{~cm}\)，高为 \(8\mathrm{~cm}\)，求圆台的体积和全面积.（\(\pi=3.14\)，结果精确到 \(0.1\mathrm{~cm}^3\) 或 \(0.1\mathrm{~cm}^2\)）
\end{enumerate}
\end{problem}

\begin{solution}
\begin{enumerate}
\item 因为 \(MN\parallel BC\parallel AD\)，三条平行线被两条截线 \(AB\)、\(CD\) 所截，所以由截比定理
\(\frac{AM}{MB}=\frac{DN}{NC}\). 两边分别加上 \(1\)，得
\(\frac{AM+MB}{MB}=\frac{DN+NC}{NC}\)，即
\(\frac{AB}{MB}=\frac{DC}{NC}\).

在三角形 \(BDA\) 中，\(M\) 在 \(BA\) 上且 \(EM\parallel AD\)，所以 \(\triangle BEM\sim\triangle BDA\)，从而
\(\frac{ME}{AD}=\frac{BM}{BA}\)，即 \(\frac{AD}{ME}=\frac{AB}{MB}\). 同理，在三角形 \(CDA\) 中，\(N\) 在 \(CD\) 上且 \(EN\parallel AD\)，所以 \(\triangle CNE\sim\triangle CDA\)，从而
\(\frac{EN}{AD}=\frac{CN}{CD}\)，即 \(\frac{AD}{EN}=\frac{DC}{NC}\).

结合 \(\frac{AB}{MB}=\frac{DC}{NC}\)，可得 \(\frac{AD}{ME}=\frac{AD}{EN}\). 因为 \(AD>0\)，所以 \(ME=EN\).

\solutionfigure{\bitmapfigure[max width=4.2cm]{img/1977/inner_mongolia_liberal/q2_solution_fig1.png}}

\item 圆台的上底半径、下底半径和高分别为
\(R_1=\frac{12}{2}\mathrm{~cm}=6\mathrm{~cm}\)，\(R_2=\frac{24}{2}\mathrm{~cm}=12\mathrm{~cm}\)，\(H=8\mathrm{~cm}\). 轴截面中，母线、高和两底半径之差构成直角三角形，因此母线长为
\(L=\sqrt{H^2+\left(R_2-R_1\right)^2}=\sqrt{8^2+6^2}=10\mathrm{~cm}\).

圆台体积公式为 \(V=\frac{1}{3}\pi H\left(R_1^2+R_1R_2+R_2^2\right)\)，代入得
\[
V=\frac{1}{3}\times3.14\times8\times\left(6^2+6\times12+12^2\right)
=672\times3.14=2110.08\mathrm{~cm}^3\approx2110.1\mathrm{~cm}^3
\]

全面积等于两个底面的面积与侧面积之和，故
\[
S=\pi\left(R_1^2+R_2^2+L\left(R_1+R_2\right)\right)
=3.14\times\left(6^2+12^2+10\times(6+12)\right)
=360\times3.14=1130.4\mathrm{~cm}^2
\]
\end{enumerate}
\end{solution}

\begin{problem}
已知：直线 \(3x+y=5\) 及 \(2x-3y+4=0\) 交于 \(A\) 点.
求：
\begin{enumerate}
\item 过 \(A\) 点且与 \(y\) 轴垂直的直线方程；
\item 过 \(A\) 点且与直线 \(2x+y+1=0\) 平行的直线方程；
\item 过 \(A\) 点且与直线 \(2x+y+1=0\) 垂直的直线方程.
\end{enumerate}
\end{problem}

\begin{solution}
\begin{enumerate}
\item 先解交点坐标. 由方程组
\[
\begin{cases}
3x+y=5,\\
2x-3y+4=0
\end{cases}
\]
将第一式化为 \(y=5-3x\)，代入第二式得 \(2x-3(5-3x)+4=0\)，即 \(11x-11=0\)，所以 \(x=1\)，从而 \(y=2\). 因此 \(A\) 点坐标为 \(\left(1,2\right)\). \(y\) 轴是竖直直线，与它垂直的直线为水平直线，所以过 \(A\) 点的所求直线为 \(y=2\).

\item 直线 \(2x+y+1=0\) 化为 \(y=-2x-1\)，斜率为 \(-2\). 平行直线斜率相同，因此过 \(A\left(1,2\right)\) 的所求直线满足点斜式
\(y-2=-2\left(x-1\right)\)，整理得 \(2x+y-4=0\).

\item 已知直线的斜率为 \(-2\)，与它垂直的直线斜率为其负倒数 \(\frac{1}{2}\). 因此过 \(A\left(1,2\right)\) 的所求直线满足
\(y-2=\frac{1}{2}\left(x-1\right)\)，化简得 \(x-2y+3=0\).
\end{enumerate}
\end{solution}

\begin{problem}
红旗拖拉机厂今年元月份生产出一批甲乙两种型号的拖拉机，其中生产乙型 \(16\text{ 台}\). 从二月份起，甲型每月增产 \(10\text{ 台}\)，而乙型按每月相同增长率逐月递增，又知二月份甲，乙两型拖拉机产量之比是 \(3:2\)，三月份甲，乙两型产量之和是 \(65\text{ 台}\)，求乙型拖拉机每月的增长率及甲型拖拉机元月份的产量.
\end{problem}

\begin{solution}
设甲型拖拉机元月份产量为 \(x\)，乙型拖拉机每月增长率为 \(y\)，则乙型拖拉机每月的增长倍数为 \(1+y\). 由于产量为正，\(1+y>0\). 二月份甲、乙两型拖拉机的产量分别为 \(x+10\)、\(16(1+y)\)，三月份分别为 \(x+20\)、\(16(1+y)^2\). 因而由题意得
\[
\begin{cases}
\frac{x+10}{16\left(1+y\right)}=\frac{3}{2},\\
16\left(1+y\right)^2+\left(x+20\right)=65.
\end{cases}
\]

由第一个方程，得
\(x=\frac{3}{2}\cdot16\left(1+y\right)-10=24\left(1+y\right)-10\). 代入第二个方程，得
\(16\left(1+y\right)^2+24\left(1+y\right)-10+20=65\)，即
\(16\left(1+y\right)^2+24\left(1+y\right)-55=0\).

令 \(t=1+y\)，则 \(t>0\)，且方程化为 \(16t^2+24t-55=0\). 因式分解得
\(\left(4t+11\right)\left(4t-5\right)=0\)，所以 \(t=-\frac{11}{4}\) 或 \(t=\frac{5}{4}\). 由于 \(t>0\)，舍去 \(t=-\frac{11}{4}\)，于是
\(1+y=\frac{5}{4}\)，即 \(y=\frac{1}{4}=25\%\). 再由 \(x=24\left(1+y\right)-10\)，得 \(x=24\times\frac{5}{4}-10=20\).

验算：二月份两型产量为 \(30\text{ 台}\) 和 \(20\text{ 台}\)，比值为 \(3:2\)；三月份两型产量为 \(40\text{ 台}\) 和 \(25\text{ 台}\)，和为 \(65\text{ 台}\)，符合题意. 故甲型拖拉机元月份产 \(20\text{ 台}\)，乙型拖拉机每月增长率为 \(25\%\).
\end{solution}
